\chapter*{前言}
\addcontentsline{toc}{chapter}{前言}

通信中的许多设计任务，共用少数几类数学结构。把一个问题称作“功率分配”还不足以选出算法：正交信道的总速率最大化可以是凸优化，干扰信道的加权和速率最大化通常是非凸优化，而给定目标 SINR 的功率最小化又可能只是线性规划。本书要帮助读者完成这一判断，并把判断落实为可求解、可验证的模型。

本书先整理从通信系统中抽象出的优化问题，再介绍所需的解法。每个主要问题都围绕四件事展开：变量控制什么，目标与约束怎样表达，什么结构决定求解难度，计算结果能够支持什么结论。注水、波束成形和干扰管理因而不再是孤立的公式，而是认识凹性、锥约束、对偶和非凸近似的具体入口。

\section*{读者与范围}

读者应掌握线性代数中的特征值、正定矩阵和矩阵分解，微积分中的梯度与 Hessian，以及概率论中的期望、条件期望和独立性。通信方面，只要求理解复基带信号、信道增益、功率和噪声；MIMO、估计和队列的必要模型在正文中就近给出。对复向量的求导最终都可以理解成对实部与虚部联合求导，无须先学习完整的复分析。

“完备”在这里指核心问题族之间的联系完整：从连续资源分配、功率控制、波束与估计，到离散分配、不确定信道和随机到达；从凸性识别、最优性条件，到一阶法、锥规划、分解以及非凸迭代。它并不要求穷举每个通信场景。读者遇到新的 RIS、边缘计算或协作传输模型时，应能先判断其中的连续变量、整数变量、干扰耦合和随机约束，再决定需要本书的哪一段工具。

本书不系统推导编码定理、随机过程的测度论基础、博弈均衡理论或强化学习算法，也不把新的系统名称各写成一章。整数全局搜索、鲁棒设计和在线控制介绍到能够建立正确模型、选择方法并理解保证为止。深入这些方向时，可沿正文引用继续阅读。需要的数学基础以通信问题为牵引补齐；长篇技术证明让位于决定算法是否适用的关键推导。

\section*{先检索，再复用}

编写前检索了作者官网、课程讲义、出版方和公开代码库。截至 2026 年 11 月 8 日，已经核实到可复用的通信优化案例库和系统性参考书；本次检索没有找到同时满足“中文短教材、问题分类在前、覆盖主要凸与非凸方法、开放改编许可、完整 LaTeX 源码”的单一成品。这个结论限定于本次核查的资源。

\begin{table}[htbp]
\centering\small
\caption{核查的主要基础与本书的使用方式}
\begin{tabularx}{\textwidth}{@{}p{0.24\textwidth}p{0.27\textwidth}X@{}}
\toprule
资源 & 已核实的开放程度 & 本书的使用方式\\
\midrule
CVX 示例库 & \texttt{examples/} 为公有领域 & 复用固定 SINR 功率控制、阵列旁瓣设计问题与建模思路\\
CVXPY 通信示例 & 示例仓库 Apache-2.0 & 在注水、最大最小 SINR 案例上补充推导、参数与数值核验\\
Boyd 凸优化书；Tse 无线通信书 & 作者提供免费 PDF；Cambridge 版权 & 核对理论与通信模型，引用原书\\
MIT OCW 6.079 讲义 & 站点采用 CC BY-NC-SA 4.0，须留意单项例外 & 作为进一步学习的公开课程资源\\
\bottomrule
\end{tabularx}
\end{table}

本书以已有案例为起点\cite{cvxexamples,cvxpyexamples}。新增工作包括：按数学结构重排通信问题；补足从物理模型到优化模型的过渡；推导关键变换与最优性条件；区分精确等价、凸松弛与局部近似；为数值实验增加独立约束检查。原始代码快照、固定版本、归属和完整许可保存在 \texttt{third\_party/} 中，具体对应关系见该目录的 \texttt{NOTICE.txt}。读者可以据此区分原有案例与本书扩展。

免费可读与开源可改编是两件事。Boyd 与 Vandenberghe 的书籍主页说明由 Cambridge 持有版权\cite{boyd2004}；Tse 与 Viswanath 的官网也说明电子版适用传统书籍的版权规则\cite{tse2005}。本书引用这些理论来源，未将原书文本或图作为可自由再分发的材料。另一套有明确开放条款的学习材料是 \href{https://ocw.mit.edu/courses/6-079-introduction-to-convex-optimization-fall-2009/pages/lecture-notes/}{MIT OCW 6.079 凸优化讲义}；其改编需遵守对应的署名、非商业和相同方式共享条件。

\section*{怎样阅读和计算}

按顺序学习时，先读第 1 章的问题目录与第 2 章的数学基础，用第 3--4 章把结构判断落实到资源、功率、波束和估计问题，再读第 5 章的数值算法。第 6 章处理干扰与非凸结构，第 7--8 章把变量离散性和随机性纳入建模。最后用实验、习题和附录检查自己是否能够从题意独立得到模型，而不是只认出某个算法名称。

已有凸优化基础的读者，可以先浏览问题目录，选择一个通信案例走完“建模—结构识别—推导—求解—复核”的全过程，再阅读非凸、离散和随机章节。初学者则应至少亲手推导注水的 KKT 条件、固定 SINR 的线性不等式以及单用户波束约束的二阶锥形式。这三条链路分别训练边际收益、耦合可行性和复数相位处理，能反复用于后面的模型。

正文中的速率按标明的带宽和对数底数解释；未给绝对尺度的实验采用归一化功率与噪声。图由明确方程或随书脚本生成，图注给出读图需要的对象、单位与设置。示例数值用于检验所述模型与算法，并不替代一般性证明。对于凸问题，最优性条件与残差提供核验依据；对于非凸问题，应同时检查可行性、目标变化和驻点条件，数条收敛曲线不能证明全局最优。

数学书写遵循用户指定的\href{https://github.com/hongjinrui/writing_skills/blob/main/SKILL.md}{科研写作规则}\cite{writing2026}：新符号首次出现时解释意义与范围，关键变换前说明目的，结论紧接适用条件与推导依据。读者遇到一个公式时，应当知道它解决当前哪一步问题，并能沿相邻文字追溯它为什么成立。
